\chapter{Relativization and natural proofs}
\label{chap:barriers}

This chapter covers the two classical barriers to separating complexity
classes. Relativization needs oracle machines and oracle worlds; the meaningful
formal endpoints are oracle existence theorems such as ``there is an oracle $A$
with $\cc{P}^A = \cc{NP}^A$'', not the meta-mathematical claim that a proof
technique relativizes. Natural proofs need Boolean-function properties measured
on truth tables, largeness as an explicit density, and pseudorandom function
security with explicit parameters. Formalized so far: deterministic
Boolean-oracle Turing machines with a dedicated query tape and one-step lookups,
their acceptance, deciding and output semantics, the conservative embedding of
ordinary machines, oracle-uniform simulation and universality interfaces (with
no concrete universal oracle machine yet), the finite counting layer of
property density, and the definition of negligible functions (without closure
lemmas). Nondeterministic oracle machines, relativized classes, the
Baker--Gill--Solovay worlds, natural properties, pseudorandom function
families and the Razborov--Rudich theorem are planned.

\section{Oracle machines}

\begin{definition}[Deterministic oracle Turing machines]
  \label{def:barriers-oracle-tm}
  \lean{Complexity.BooleanOracle, Complexity.BooleanOracle.Decides,
    Complexity.Tape.oracleQuery, Complexity.OracleCfg, Complexity.OracleCfg.init,
    Complexity.OracleLocalTransition, Complexity.OracleTM,
    Complexity.OracleTM.step, Complexity.OracleTM.stepRel,
    Complexity.OracleTM.reachesIn}
  \leanok
  \uses{def:models-tape, def:models-cfg, def:tm, def:language}
  A Boolean oracle is a function $O : \{0,1\}^* \to \{0,1\}$; it decides a
  language $A$ if $O(w) = 1 \iff w \in A$ for every $w$. An oracle TM with $k$
  work tapes has a finite state type with decidable equality, a start state
  and a halt state; the input, $k$ work and output tapes of an ordinary TM plus
  a dedicated read-write query tape; a partial map sending each query state to
  a pair (yes-successor, no-successor); and a local transition that reads the
  symbols under all heads (the query head included), writes symbols of
  $\Gamma_w$ on the query, work and output tapes, and moves every head, moving
  right every head that reads $\triangleright$. The initial configuration on
  $x$ is in the start state with $x$ on the input tape and every other tape
  empty. One step relative to $O$ is defined as follows. The halt state has no
  successor. In a non-halted query state the machine reads the query word,
  formed by the cells strictly between the left marker and the query head in
  left-to-right order (a cell reads $1$ only if it holds $1$; $0$, blank and
  $\triangleright$ read as $0$), and enters the yes-successor if $O$ answers
  $1$ on it and the no-successor otherwise. The lookup costs exactly one step
  and changes nothing but the state; in particular the query tape is not
  erased, and writing the query costs ordinary local steps. Every other
  non-halted state applies the local transition. The relation
  $c \to_O^t c'$ holds when $c'$ is reached from $c$ in exactly $t$ steps
  relative to $O$.
\end{definition}

\begin{definition}[Oracle acceptance, deciding and output]
  \label{def:barriers-oracle-decides}
  \lean{Complexity.OracleTM.Accepts, Complexity.OracleTM.AcceptsInTime,
    Complexity.OracleTM.DecidesInTime, Complexity.OracleTM.HaltsInTime,
    Complexity.OracleTM.Halts, Complexity.OracleTM.ProducesInTime,
    Complexity.OracleTM.Produces}
  \leanok
  \uses{def:barriers-oracle-tm, def:models-tape, def:language}
  Fix an oracle $O$ and an oracle TM $M$; all runs are relative to $O$, and
  every step, local or lookup, counts toward a budget. $M$ accepts $x$ (within
  $t$ steps) if the run on $x$ reaches the halt state (after at most $t$
  steps) with output cell $1$ equal to $1$. $M$ decides $L$ in time
  $T : \mathbb{N} \to \mathbb{N}$ if on every input $x$ it reaches the halt
  state within $T(|x|)$ steps, with output cell $1$ equal to $1$ if $x \in L$
  and equal to $0$ if $x \notin L$. $M$ halts on a program $p$ within $t$
  steps if the run on $p$ reaches the halt state after at most $t$ steps, and
  produces $y$ on $p$ within $t$ steps if moreover the halted output tape has
  output $y$. $M$ halts on $p$ (produces $y$ on $p$) if it does so within some
  finite budget. These have the same shape as the corresponding definitions for
  ordinary machines.
\end{definition}

\begin{lemma}[Clock monotonicity for oracle machines]
  \label{lem:barriers-oracle-clock-mono}
  \lean{Complexity.OracleTM.HaltsInTime.mono,
    Complexity.OracleTM.halts_of_haltsInTime,
    Complexity.OracleTM.ProducesInTime.mono,
    Complexity.OracleTM.produces_of_producesInTime}
  \leanok
  \uses{def:barriers-oracle-decides}
  For every oracle, halting and producing $y$ within $t$ steps imply the same
  within any $t' \ge t$ steps, and imply eventual halting and eventual
  production of $y$, respectively.
\end{lemma}
\begin{proof}
  \leanok
  The same run witnesses the larger budget, and the eventual notions are
  defined by existentially quantifying the budget.
\end{proof}

\begin{lemma}[Oracle execution is deterministic]
  \label{lem:barriers-oracle-deterministic}
  \lean{Complexity.OracleTM.step_eq_none_iff_halted,
    Complexity.OracleTM.stepRel_functional,
    Complexity.OracleTM.reachesIn_functional,
    Complexity.OracleTM.step_query_true, Complexity.OracleTM.step_query_false,
    Complexity.OracleTM.reachesIn_one_query_true,
    Complexity.OracleTM.reachesIn_one_query_false,
    Complexity.Tape.length_oracleQuery, Complexity.OracleTM.reachesIn_map}
  \leanok
  \uses{def:barriers-oracle-tm}
  Fix an oracle. A step is undefined exactly at the halt state; each
  configuration has at most one successor and each exact-time run has a unique
  final configuration. In a non-halted query state, a true answer is one step
  to the yes-successor and a false answer is one step to the no-successor,
  with all tapes unchanged. The query word of a tape whose head is at cell $h$
  has exactly $h - 1$ bits (none when $h = 0$). A map from oracle
  configurations to configurations of an ordinary machine that sends every
  oracle step to an ordinary step sends every exact-time oracle run to an
  ordinary run with the same number of steps.
\end{lemma}
\begin{proof}
  \leanok
  Unfold the step function; the last claim is induction on the run.
\end{proof}

\begin{definition}[Ordinary machines as oracle machines]
  \label{def:barriers-to-oracle}
  \lean{Complexity.TM.toOracleTM, Complexity.OracleCfg.erase}
  \leanok
  \uses{def:tm, def:models-cfg, def:barriers-oracle-tm}
  Every TM with $k$ work tapes is an oracle TM with $k$ work tapes, the same
  states, start and halt state, and no query states. Its local transition
  applies the source transition to the input, work and output tapes and keeps
  the query tape inert: it writes only blanks there, and moves the query head
  right when it reads $\triangleright$ and leaves it in place otherwise.
  Erasing the query tape maps an oracle configuration to the ordinary
  configuration with the same state and input, work and output tapes.
\end{definition}

\begin{theorem}[The embedding is conservative for every oracle]
  \label{thm:barriers-to-oracle-conservative}
  \lean{Complexity.TM.toOracleTM_queryTransition,
    Complexity.TM.toOracleTM_step_oracle_independent,
    Complexity.OracleCfg.erase_init,
    Complexity.TM.erase_toOracleTM_step,
    Complexity.TM.exists_toOracleTM_step_of_step_erase,
    Complexity.TM.erase_toOracleTM_reachesIn,
    Complexity.TM.exists_toOracleTM_reachesIn_of_reachesIn_erase,
    Complexity.TM.toOracleTM_decidesInTime_iff,
    Complexity.TM.toOracleTM_producesInTime_iff}
  \leanok
  \uses{def:barriers-to-oracle, def:barriers-oracle-decides, def:tm-decides,
    def:produces}
  Let $M$ be a TM and $M'$ its embedding. $M'$ has no query states and each of
  its steps is the same for every oracle. Erasing the initial oracle
  configuration on $x$ gives the ordinary initial configuration on $x$. For
  every oracle and every oracle configuration $c$, the step of $M'$ from $c$,
  erased, equals the step of $M$ from the erasure of $c$ (both undefined or
  both defined and equal), and every step of $M$ from the erasure of $c$ lifts
  to a step of $M'$ from $c$ whose result erases to it. Hence every exact-time
  run of $M'$ erases to a run of $M$ with the same number of steps, and every
  exact-time run of $M$
  from the erasure of $c$ lifts to a run of $M'$ from $c$ of the same length
  whose final configuration erases to the final configuration of $M$.
  Consequently, for every oracle $O$, language $L$ and bound
  $T : \mathbb{N} \to \mathbb{N}$, $M'$ decides $L$ in time $T$ relative to $O$
  if and only if $M$ decides $L$ in time $T$; and for every program $p$, output
  $y$ and budget $t$, $M'$ produces $y$ on $p$ within $t$ steps relative to $O$
  if and only if $M$ does.
\end{theorem}
\begin{proof}
  \leanok
  Step-by-step correspondence under erasure, in both directions: erase an
  oracle run to get the source run, and lift a source run from the erased
  initial configuration, reading the verdict or output off the erased final
  configuration.
\end{proof}

\begin{definition}[Oracle universality]
  \label{def:barriers-oracle-universal}
  \lean{Complexity.OracleTM.Simulates, Complexity.OracleTM.SimulatesInTime,
    Complexity.OracleTM.IsUniversal,
    Complexity.OracleTM.IsEfficientlyUniversalFor,
    Complexity.OracleTM.IsEfficientlyUniversal}
  \leanok
  \uses{def:barriers-oracle-decides, def:computable, def:universal,
    def:models-simulates-in-time}
  An oracle machine $U$ simulates an oracle machine $M$ under a compiler
  $c : \{0,1\}^* \to \{0,1\}^*$ if for every oracle $O$ and program $p$: $U$
  halts on $c(p)$ relative to $O$ if and only if $M$ halts on $p$ relative to
  $O$, and for every $y$, $U$ produces $y$ on $c(p)$ relative to $O$ if and
  only if $M$ produces $y$ on $p$ relative to $O$. $U$ simulates $M$ under $c$
  in time $\tau : \{0,1\}^* \times \mathbb{N} \to \mathbb{N}$ if for every
  oracle $O$, whenever $M$ produces $y$ on $p$ within $t$ steps relative to
  $O$, $U$ produces $y$ on $c(p)$ within $\tau(p, t)$ steps relative to $O$;
  unlike the ordinary timed simulation, bounded halting alone is not required
  to transfer. $U$ is universal if for every $k$ and every oracle machine $M$
  with $k$ work tapes there is a compiler $c$, computable by an ordinary
  deterministic machine, such that $U$ simulates $M$ under $c$. For a policy
  of admissible clocks, $U$ is efficiently universal for the policy if moreover
  $c$ has additive overhead ($|c(p)| \le |p| + C$) and $U$ simulates $M$ under
  $c$ in some admissible clock $\tau$; it is efficiently universal if this
  holds for the polynomial policy, which admits $\tau$ when
  $\tau(p, t) \le a (|p| + t + 1)^e$ for some constants $a, e$. In every case
  the compiler, constant and clock depend on $M$ but are chosen once for all
  oracles.
\end{definition}

\begin{lemma}[Basic facts about oracle simulation]
  \label{lem:barriers-oracle-universal-basic}
  \lean{Complexity.OracleTM.Simulates.refl, Complexity.OracleTM.Simulates.comp,
    Complexity.OracleTM.SimulatesInTime.refl,
    Complexity.OracleTM.SimulatesInTime.comp,
    Complexity.OracleTM.IsEfficientlyUniversalFor.isUniversal,
    Complexity.OracleTM.IsEfficientlyUniversal.isUniversal,
    Complexity.OracleTM.IsEfficientlyUniversalFor.mono}
  \leanok
  \uses{def:barriers-oracle-universal}
  Every oracle machine simulates itself under the identity compiler, in time
  $(p, t) \mapsto t$. If $U_1$ simulates $U_2$ under $c_1$ and $U_2$ simulates
  $M$ under $c_2$, then $U_1$ simulates $M$ under $c_1 \circ c_2$, and
  likewise for timed simulations with clock
  $(p, t) \mapsto \tau_1(c_2(p), \tau_2(p, t))$. Efficient oracle universality
  for any clock policy, in particular the polynomial one, implies oracle
  universality, and enlarging the policy preserves efficient universality.
\end{lemma}
\begin{proof}
  \leanok
  Compose the compilers and unfold the definitions.
\end{proof}

\begin{theorem}[A universal oracle machine]
  \label{thm:barriers-universal-oracle-machine}
  \uses{def:barriers-oracle-universal}
  There is an oracle TM that is efficiently universal (for the polynomial
  policy). At present oracle simulation and universality are only interfaces:
  a timed oracle simulation with additive overhead and a polynomial clock is a
  hypothesis of the oracle Kolmogorov-complexity transfer theorem
  (Lemma~\ref{lem:avg-oracle-kolmogorov}(c)), and
  efficient oracle universality is a hypothesis of the Nisan--Wigderson
  reconstruction; no universal oracle machine has been constructed.
\end{theorem}
\begin{proof}
  \uses{thm:utm}
  Extend the ordinary universal machine with a query tape that it maintains on
  behalf of the simulated machine, forwarding each simulated lookup as one real
  lookup.
\end{proof}

\section{Relativized classes and the Baker--Gill--Solovay worlds}

\begin{definition}[Nondeterministic oracle machines]
  \label{def:barriers-oracle-ntm}
  \uses{def:barriers-oracle-tm, def:ntm}
  An oracle NTM has two local transition functions selected by a choice bit and
  the same query mechanism and cost as an oracle TM, with execution along a
  fixed choice sequence and existential acceptance within a time bound.
\end{definition}

\begin{definition}[Relativized classes]
  \label{def:barriers-relativized-classes}
  \uses{def:barriers-oracle-decides, def:barriers-oracle-ntm, def:P, def:NP}
  \notready
  For a language $A$, $\cc{P}^A$ is the class of languages decided in
  polynomial time by an oracle TM relative to an oracle deciding $A$, and
  $\cc{NP}^A$ is the nondeterministic analogue. Query length is charged through
  the time needed to write the query.
\end{definition}

\begin{proposition}[Relativizing basics]
  \label{prop:barriers-relativizing-basics}
  \uses{def:barriers-relativized-classes}
  \notready
  For every $A$: $\cc{P} \subseteq \cc{P}^A \subseteq \cc{NP}^A$ and
  $\cc{NP} \subseteq \cc{NP}^A$; if $A \in \cc{P}$ then $\cc{P}^A = \cc{P}$.
\end{proposition}
\begin{proof}
  \uses{thm:barriers-to-oracle-conservative}
  The first inclusion is the conservative embedding. The collapse for
  $A \in \cc{P}$ replaces each lookup by a call to a decider, which needs
  machine-level subroutine inlining.
\end{proof}

\begin{definition}[Clocked enumeration of oracle machines]
  \label{def:barriers-clocked-enumeration}
  \uses{def:barriers-oracle-universal}
  \notready
  An enumeration $i \mapsto (M_i, p_i)$ of oracle TMs paired with polynomial
  clocks such that every polynomial-time oracle language, relative to every
  oracle, is decided by some clocked $M_i$, and the clocked machines are
  simulated uniformly in the oracle.
\end{definition}

\begin{lemma}[The unary test language is in $\cc{NP}^B$]
  \label{lem:barriers-unary-test-language}
  \uses{def:barriers-relativized-classes}
  \notready
  For every $B$, the language $U_B = \{1^n : B \text{ contains a string of
  length } n\}$ is in $\cc{NP}^B$.
\end{lemma}
\begin{proof}
  Guess a string of length $n$, write it on the query tape and query it.
\end{proof}

\begin{theorem}[A separating oracle]
  \label{thm:barriers-separating-oracle}
  \uses{def:barriers-relativized-classes}
  \notready
  There is an oracle $B$ with $\cc{P}^B \neq \cc{NP}^B$.
\end{theorem}
\begin{proof}
  \uses{lem:barriers-unary-test-language, def:barriers-clocked-enumeration}
  Build $B$ in stages so that the $i$-th clocked machine errs on $U_B$ at a
  fresh length $n$ large enough that its clock $p_i(n)$ is below $2^n$; such a
  run queries fewer than $2^n$ strings of length $n$, and later stages must not
  change answers earlier stages relied on.
\end{proof}

\begin{lemma}[A PSPACE-complete oracle collapses both classes]
  \label{lem:barriers-pspace-oracle-collapse}
  \uses{def:barriers-relativized-classes, def:PSPACE}
  \notready
  If $A$ is PSPACE-complete under polynomial-time many-one reductions then
  $\cc{P}^A = \cc{NP}^A = \cc{PSPACE}$.
\end{lemma}
\begin{proof}
  \uses{thm:space-tqbf-complete, thm:space-savitch}
  $\cc{PSPACE} \subseteq \cc{P}^A$ by reduction to $A$. $\cc{NP}^A \subseteq
  \cc{NPSPACE} = \cc{PSPACE}$ needs the closure theorem that a polynomial-space
  machine can answer the queries itself; this closure is the part that does
  not follow from completeness alone.
\end{proof}

\begin{theorem}[An equalizing oracle]
  \label{thm:barriers-equality-oracle}
  \uses{def:barriers-relativized-classes}
  \notready
  There is an oracle $A$ with $\cc{P}^A = \cc{NP}^A$.
\end{theorem}
\begin{proof}
  \uses{lem:barriers-pspace-oracle-collapse, thm:space-tqbf-complete}
  Take $A = \cc{TQBF}$.
\end{proof}

\begin{corollary}[Baker--Gill--Solovay]
  \label{cor:barriers-bgs}
  \uses{def:barriers-relativized-classes}
  \notready
  There are oracles $A$ and $B$ with $\cc{P}^A = \cc{NP}^A$ and
  $\cc{P}^B \neq \cc{NP}^B$.
\end{corollary}
\begin{proof}
  \uses{thm:barriers-equality-oracle, thm:barriers-separating-oracle}
  Combine the two worlds.
\end{proof}

\section{Natural proofs}

\begin{lemma}[Counting Boolean functions]
  \label{lem:barriers-card-boolfunc}
  \lean{Complexity.card_boolFunc}
  \leanok
  There are exactly $2^{2^n}$ Boolean functions on $n$ bits, that is,
  $2^{2^n}$ truth tables of length $2^n$.
\end{lemma}
\begin{proof}
  \leanok
  Cardinality of a function type.
\end{proof}

\begin{definition}[Property density]
  \label{def:barriers-density}
  \lean{Complexity.density}
  \leanok
  \uses{lem:barriers-card-boolfunc}
  For a set $P$ of Boolean functions on $n$ bits (necessarily finite), its
  density is the exact rational $|P| / 2^{2^n}$, the fraction of all Boolean
  functions on $n$ bits that lie in $P$.
\end{definition}

\begin{lemma}[Density laws]
  \label{lem:barriers-density-laws}
  \lean{Complexity.card_le_card_boolFunc, Complexity.density_nonneg,
    Complexity.density_le_one, Complexity.density_compl,
    Complexity.density_union_le}
  \leanok
  \uses{def:barriers-density, lem:barriers-card-boolfunc}
  For all sets $P, Q$ of Boolean functions on $n$ bits: $|P| \le 2^{2^n}$,
  $0 \le \mathrm{density}(P) \le 1$,
  $\mathrm{density}(P^c) = 1 - \mathrm{density}(P)$ with $P^c$ the complement
  among all Boolean functions on $n$ bits, and
  $\mathrm{density}(P \cup Q) \le \mathrm{density}(P) + \mathrm{density}(Q)$.
\end{lemma}
\begin{proof}
  \leanok
  Finite cardinality arithmetic.
\end{proof}

\begin{definition}[Negligible functions]
  \label{def:barriers-negligible}
  \lean{Complexity.Negligible}
  \leanok
  $f : \mathbb{N} \to \mathbb{R}$ is negligible if for every $c \in \mathbb{N}$
  there is $N > 0$ with $|f(n)| < 1 / n^c$ for all $n \ge N$.
\end{definition}

\begin{lemma}[Closure of negligible functions]
  \label{lem:barriers-negligible-closure}
  \uses{def:barriers-negligible}
  Sums of negligible functions and products of a negligible function with a
  polynomial are negligible; a function bounded in absolute value by a
  negligible function is negligible.
\end{lemma}

\begin{definition}[Natural properties]
  \label{def:barriers-natural-property}
  \uses{def:barriers-density, def:SIZE, def:PPoly}
  A family of properties $P_n$ of $n$-bit Boolean functions is large with
  density $\delta(n)$ if $\mathrm{density}(P_n) \ge \delta(n)$; constructive if
  membership of a truth table of length $N = 2^n$ is decidable in time
  polynomial in $N$ (or by circuits of size polynomial in $N$); and useful
  against $\cc{SIZE}(s)$ if, for all sufficiently large $n$, no function in
  $P_n$ has circuits of size $s(n)$. The quantifier over lengths and the scale
  of constructivity are part of the definition.
\end{definition}

\begin{definition}[Pseudorandom function families]
  \label{def:barriers-prf}
  \uses{def:barriers-negligible, def:PPoly}
  A keyed family $f_k : \{0,1\}^n \to \{0,1\}$ computable in polynomial time,
  and its security game against a nonuniform class of distinguishers with oracle
  access (or with the full truth table): the advantage in telling $f_k$ for a
  random key from a uniformly random function, required to be negligible, or
  bounded by an explicit function of the distinguisher's size.
\end{definition}

\begin{lemma}[A large property excluding a range is a distinguisher]
  \label{lem:barriers-range-distinguisher}
  \uses{def:barriers-density}
  Let $P$ be a set of Boolean functions on $n$ bits with density at least
  $\delta$, and let a family $\{f_k\}$ have no member in $P$. Then the test
  ``the truth table lies in $P$'' accepts a uniformly random function with
  probability at least $\delta$ and accepts $f_k$ with probability $0$, so its
  distinguishing advantage is at least $\delta$. No circuit assumption is used.
\end{lemma}
\begin{proof}
  Count. The finite statistical-test layer of the metacomplexity development
  already proves the analogous statement for dense sets of high-complexity
  strings against low-complexity generators
  (Lemma~\ref{lem:avg-test-hybrid}(a)), but not this Boolean-function form.
\end{proof}

\begin{theorem}[Razborov--Rudich, parameterized conditional form]
  \label{thm:barriers-razborov-rudich}
  \uses{def:barriers-natural-property, def:barriers-prf}
  \notready
  Suppose a pseudorandom function family computable by circuits of size
  $s(n)$ has advantage less than $\delta(n)$ against every truth-table
  distinguisher of size polynomial in $2^n$. Then no property that is large
  with density $\delta$ and constructive at that size is useful against
  $\cc{SIZE}(s)$.
\end{theorem}
\begin{proof}
  \uses{lem:barriers-range-distinguisher}
  Usefulness puts every $f_k$ outside the property; largeness and
  constructivity turn membership into a distinguisher whose advantage is the
  density, contradicting security.
\end{proof}

\begin{corollary}[The natural-proofs barrier]
  \label{cor:barriers-natural-proofs}
  \uses{def:barriers-natural-property, def:barriers-prf, def:PPoly}
  \notready
  If for some $\varepsilon > 0$ there is a pseudorandom function family in
  $\cc{P/poly}$ against which no distinguisher of size $2^{n^{\varepsilon}}$
  has advantage $2^{-n^{\varepsilon}}$ or more, then no
  $\cc{P/poly}$-constructive property of density at least $2^{-O(n)}$ is
  useful against $\cc{P/poly}$. The pseudorandomness hypothesis is an explicit
  assumption, not a proved fact.
\end{corollary}
\begin{proof}
  \uses{thm:barriers-razborov-rudich, lem:barriers-negligible-closure}
  Instantiate the parameters: a truth table of length $N = 2^n$ is tested by
  circuits of size $\mathrm{poly}(N)$, which the strong security hypothesis
  covers after rescaling $n$.
\end{proof}
